% STATUS: written, 20 October 2026. Needs a figure reference in paragraph 4
% once Figure 1 exists.

\section{Introduction}

A distributed database splits its data into pieces and spreads those pieces
across machines in different places. The pieces are usually called shards, or
fragments, and the two decisions that matter are where the boundaries between
them fall and which machine holds each one. Both decisions are made once, when
the database is designed. The boundaries are fixed by a schema, and the
placement is decided by a function of the key, such as a hash, which maps a
record to a machine without ever considering where the requests for that record
come from.

That arrangement works whilst the traffic stays where the designer expected it.
In a service with users on several continents, served by a database spread
across regions in the manner of Spanner~\cite{corbett2013spanner} or
CockroachDB~\cite{taft2020cockroachdb}, it does not. The geographic
centre of the workload moves across the day as one region wakes and another
sleeps, and it moves abruptly when something local draws attention to one
region. A placement chosen for the morning is wrong by the evening. Every
request that has to cross an ocean to reach its data pays the round-trip time
of that ocean, and a transaction touching several shards held in several
regions pays it repeatedly, because the protocol that keeps such a transaction
atomic requires more than one round of agreement between the machines
involved. The result is that latency degrades steadily as the mismatch between
placement and traffic grows.

The obvious response is to move the data when the traffic moves. A simple rule
does this: each decision window, find the shard whose placement costs the most
and move it towards the region reading it. I implement that rule as a baseline
and call it \emph{follow traffic}. The question this paper asks is whether
anything is gained by replacing it with a policy learned from experience, and
the answer turns out to depend entirely on how fast the workload moves. When
drift is slow, the simple rule is within half a millisecond of the best
placement that moving shards can achieve, and there is nothing worth learning.
When drift is fast, the same rule falls 57\,ms short of that bound, which is
41\% of its own latency, because it reacts one window behind a target that has
already moved again. The room for a learned policy is therefore largest exactly
where the problem is hardest, and that is a measured observation rather than an
assumption.

This paper reports the following. I build a simulator of a six-region cluster
with a workload whose geographic centre drifts on a daily cycle and spikes
locally, and I define the placement problem on it as a Markov decision process
with a flat action space of 261 declarative actions. I establish four reference
points against which any learned policy must be judged: static hash
partitioning, which is what production systems do today; a random control; the
follow-traffic heuristic; and a per-window bound on what lease movement alone
can achieve. I measure how the gap between the heuristic and that bound widens
with workload volatility. I verify on a real six-node CockroachDB cluster, with
wide-area latency injected between the nodes, that placement changes
transaction latency by far more than the measurement noise, which is the
precondition for any of this being learnable outside simulation. Finally I
train a Deep Q-Network on the simulator and report, honestly, that it does not
yet match the heuristic, together with the evidence about why.

% TODO once Figure 1 exists: close the paragraph above with a pointer to it.
